% GENERATED FILE. Edit canonical lessons, then run npm run generate:books.
% canonical-source-hash: fnv1a64:cf2a8dd578f13505
% canonical-lessons: MR-C10-paani, MR-C10-chaha, MR-C10-dudh, MR-C10-bhakari

\chapter{Water, Tea, Milk, and Bhakri}
\label{ch:mr-request-food-drink}
\hlchaptermodality{10}

\begin{chapteropening}
\textbf{By the end of this chapter:} \emph{I can politely ask for water, tea, milk, or bhakri in Marathi, name which of the four is a Chinese loan and which never touched Sanskrit at all, and say why all three of Marathi's genders already show up in one short chapter.}

The chapter closes on \mr{भाकरी}, the request that breaks every pattern the first three set: the reader produces all four requests with \mr{कृपया}, sorts \mr{पाणी} and \mr{दूध} (inherited from Sanskrit) against \mr{चहा} (a Chinese loan by the overland route) against \mr{भाकरी} (native to Marathi's own ground, sharing Sanskrit with none of them), and names the three genders the four words already cover.
\end{chapteropening}

\section[pāṇī]{\mr{पाणी} (pāṇī) — the first thing you'll actually ask for}

\label{lesson:MR-C10-paani}



\begin{quote}
\textbf{\mr{आवडणे}} told you what you like; \textbf{\mr{प्रेम} \mr{करणे}} what you love. This chapter asks for something smaller and more frequent than either: a drink.
\end{quote}

\subsection*{You'll want to know: \mr{पाणी}}
\begin{quote}
\textbf{\mr{पाणी}} — \emph{pāṇī} — \textbf{water}
\end{quote}

\subsection*{The letters in this word}
Nothing new here. \textbf{\mr{प}} \emph{pa} you have from \emph{punhā}; \textbf{\mr{ण}} you have from every \textbf{-\mr{णे}} verb so far; \textbf{\mr{ी}} is the familiar long \emph{ī}. \textbf{\mr{पा}-\mr{णी}}.

\begin{grammarlens}[title={Grammar Lens: asking politely}]
\begin{quote}
\textbf{\mr{पाणी}, \mr{कृपया}.} — \emph{pāṇī, kṛpayā.} — \textbf{Water, please.}
\end{quote}

\textbf{\mr{कृपया}} \emph{kṛpayā} is Sanskrit's own instrumental of \textbf{\mr{कृपा}} \emph{kṛpā}, \textquotedblleft{}mercy, kindness\textquotedblright{} — literally \textquotedblleft{}by kindness, out of kindness.\textquotedblright{} Name the thing you want, add \textbf{\mr{कृपया}}, and the request is complete: no verb for \textquotedblleft{}give\textquotedblright{} required. This is the pattern the whole chapter reuses, and the first time this track has built a \textquotedblleft{}please\textquotedblright{} of any kind — a spine node this book had left empty until now.
\end{grammarlens}

\begin{cousinweb}
\textbf{\mr{पाणी}} continues Maharastri Prakrit \textbf{pāṇīa}, from Sanskrit \textbf{\mr{पानीय}} \emph{pānīya}, \textquotedblleft{}that which is to be drunk,\textquotedblright{} built on \textbf{\mr{पा}} \emph{pā-}, \textquotedblleft{}to drink.\textquotedblright{} That root is Indo-European \emph{\textbf{peh\textsubscript{3}(i)-}}, and English keeps two disguised pieces of it: \textbf{potion} and \textbf{poison} both come through Latin \emph{potare}, \textquotedblleft{}to drink\textquotedblright{} — a poison, by origin, is just a drink gone wrong. Greek turned the same idea into \emph{sympotēs}, \textquotedblleft{}one who drinks together,\textquotedblright{} which English keeps as \textbf{symposium}.
\end{cousinweb}

\subsection*{Your turn}
\begin{itemize}
\raggedright
  \item \emph{Say it:} \textbf{pāṇī} — water
  \item \emph{Say it:} \textbf{pāṇī, kṛpayā.} — water, please
  \item \emph{Run through:} \textbf{malā marāṭhī āvaḍte}, then \textbf{pāṇī, kṛpayā} — liking and requesting, side by side
\end{itemize}

\begin{tcolorbox}[breakable,colback=teal!4,colframe=teal!35!black,title={Before you move on}]
Ask for water politely. (\textbf{Pāṇī, kṛpayā.}) What does \textbf{\mr{कृपया}} literally mean? (\textquotedblleft{}By kindness.\textquotedblright{}) Which two English words for a drink, and one for a whole gathering of them, share \textbf{pāṇī}'s root? (\textbf{Potion}, \textbf{poison}, \textbf{symposium}.)

Sources: \href{https://en.wiktionary.org/wiki/\%E0\%A4\%AA\%E0\%A4\%BE\%E0\%A4\%A3\%E0\%A5\%80}{Wiktionary: \mr{पाणी}}; \href{https://en.wiktionary.org/wiki/\%E0\%A4\%95\%E0\%A5\%83\%E0\%A4\%AA\%E0\%A4\%AF\%E0\%A4\%BE}{Wiktionary: \mr{कृपया}}.
\end{tcolorbox}
\section[chahā]{\mr{चहा} (chahā) — a word that took the overland road}

\label{lesson:MR-C10-chaha}



\begin{quote}
\textbf{\mr{पाणी}, \mr{कृपया}} (\emph{pāṇī, kṛpayā}) — water, please. Say it once more, then ask for something that is not water at all.
\end{quote}

\subsection*{You'll want to know: \mr{चहा}}
\begin{quote}
\textbf{\mr{चहा}} — \emph{chahā} — \textbf{tea}
\end{quote}

\begin{quote}
\textbf{\mr{चहा}, \mr{कृपया}.} — \emph{chahā, kṛpayā.} — \textbf{Tea, please.}
\end{quote}

\subsection*{The letters in this word}
Nothing new. \textbf{\mr{च}} you already say near \emph{ts} from Chapter 6's \textbf{\mr{चार}}; \textbf{\mr{ह}} is \textbf{\mr{नमस्कार}}'s; \textbf{\mr{ा}} the familiar long \emph{ā}. \emph{Cha-hā}.

\begin{cousinweb}
Every other word in this chapter reaches back to Sanskrit. \textbf{\mr{चहा}} does not — it is a loanword from Mandarin \textbf{cha}, and it reached Marathi by the \textbf{overland} route, through Persian and Turkic trade, the same road that gave Hindi \textbf{\mr{चाय}} \emph{chāy}, Bengali \emph{cha}, and Persian \emph{chāy} their words for the same drink. A second branch of the same plant travelled by \textbf{sea}: Dutch traders carried the Hokkien word \textbf{te} to Western Europe, which is why English says \textbf{tea}, not \emph{cha}. One plant, two words, two trade routes.

\textbf{\mr{चहा}} is also this chapter's first \textbf{masculine} noun, against \textbf{\mr{पाणी}}'s neuter — a reminder that Marathi's three genders do not sort words by where they came from.
\end{cousinweb}

\subsection*{Your turn}
\begin{itemize}
\raggedright
  \item \emph{Say it:} \textbf{chahā, kṛpayā} — tea, please
  \item \emph{Name:} the two trade routes — overland to \emph{chāy/chahā}, by sea to \emph{tea}
  \item \emph{Say it:} both requests — \textbf{pāṇī, kṛpayā … chahā, kṛpayā}
\end{itemize}

\begin{tcolorbox}[breakable,colback=teal!4,colframe=teal!35!black,title={Before you move on}]
Ask for tea politely. (\textbf{Chahā, kṛpayā.}) Which language does \textbf{\mr{चहा}} ultimately come from? (\textbf{Mandarin Chinese}, \emph{cha}.) Which English word for the same plant took the sea route instead? (\textbf{Tea}.)

Sources: \href{https://en.wiktionary.org/wiki/\%E0\%A4\%9A\%E0\%A4\%B9\%E0\%A4\%BE}{Wiktionary: \mr{चहा}}.
\end{tcolorbox}
\section[dūdh]{\mr{दूध} (dūdh) — back to Sanskrit, and a surprising English cousin}

\label{lesson:MR-C10-dudh}



\begin{quote}
\textbf{\mr{चहा}} broke the pattern — a loanword with no Sanskrit behind it. This next word rejoins \textbf{\mr{पाणी}}'s family: inherited, traceable, Indo-European.
\end{quote}

\subsection*{You'll want to know: \mr{दूध}}
\begin{quote}
\textbf{\mr{दूध}} — \emph{dūdh} — \textbf{milk}
\end{quote}

\begin{quote}
\textbf{\mr{दूध}, \mr{कृपया}.} — \emph{dūdh, kṛpayā.} — \textbf{Milk, please.}
\end{quote}

\subsection*{The letters in this word}
One new mark: \textbf{\mr{ू}}, the \textbf{long \emph{ū}} vowel sign — a longer hook under the consonant than any you've written so far. \textbf{\mr{द}} you have from \emph{dhanyavād} and \emph{madat}; \textbf{\mr{ध}} from \emph{dhanyavād} too. \emph{Dū-dh}.

\begin{cousinweb}
\textbf{\mr{दूध}} continues Maharastri Prakrit \textbf{duddha}, from Sanskrit \textbf{\mr{दुग्ध}} \emph{dugdha}, \textquotedblleft{}milked, milk\textquotedblright{} — literally the past participle of \textbf{\mr{दुह्}} \emph{duh-}, \textquotedblleft{}to milk.\textquotedblright{} That root traces to Indo-European \emph{\textbf{*d\textsuperscript{h}ewg\textsuperscript{h}-}}, \textquotedblleft{}to be fit, to be of use, to produce\textquotedblright{} — milking as \emph{producing something useful}. Germanic kept the same root for a different sense of \textquotedblleft{}useful, capable\textquotedblright{}: English \textbf{doughty}, \textquotedblleft{}brave and able,\textquotedblright{} is \emph{dūhtig} in Old English, built on this exact root. Milk and doughtiness, four thousand years apart, share one ancestor.
\end{cousinweb}

\subsection*{Your turn}
\begin{itemize}
\raggedright
  \item \emph{Say it:} \textbf{dūdh, kṛpayā} — milk, please
  \item \emph{Say it:} what \emph{dugdha} literally was — \textquotedblleft{}milked\textquotedblright{}
  \item \emph{Run through:} all three requests — \textbf{pāṇī, chahā, dūdh} — kṛpayā
\end{itemize}

\begin{tcolorbox}[breakable,colback=teal!4,colframe=teal!35!black,title={Before you move on}]
Ask for milk politely. (\textbf{Dūdh, kṛpayā.}) What English word for \textquotedblleft{}brave and capable\textquotedblright{} shares \textbf{\mr{दूध}}'s root? (\textbf{Doughty}.)

Sources: \href{https://en.wiktionary.org/wiki/\%E0\%A4\%A6\%E0\%A5\%82\%E0\%A4\%A7}{Wiktionary: \mr{दूध}}.
\end{tcolorbox}
\section[bhākarī]{\mr{भाकरी} (bhākarī) — the request that breaks the pattern}

\label{lesson:MR-C10-bhakari}



\begin{quote}
Three requests so far: \textbf{pāṇī}, \textbf{chahā}, \textbf{dūdh}. Two inherited from Sanskrit, one borrowed from Chinese. The fourth belongs to neither story.
\end{quote}

\subsection*{You'll want to know: \mr{भाकरी}}
\begin{quote}
\textbf{\mr{भाकरी}} — \emph{bhākarī} — \textbf{a round flatbread}, the everyday Maharashtrian bread — of jowar, bajra, or rice flour, not wheat
\end{quote}

\begin{quote}
\textbf{\mr{भाकरी}, \mr{कृपया}.} — \emph{bhākarī, kṛpayā.} — \textbf{Bhakri, please.}
\end{quote}

\subsection*{The letters in this word}
Nothing new. \textbf{\mr{भ}} from \emph{bheṭū}, \textbf{\mr{क}} from everywhere, \textbf{\mr{र}} from \emph{marāṭhī}, \textbf{\mr{ी}} the familiar long \emph{ī}. \emph{Bhā-ka-rī}.

\begin{cousinweb}
\textbf{\mr{भाकरी}} is inherited from Old Marathi \textbf{bhākarī}, from Proto-Indo-Aryan \emph{\textbf{*b\textsuperscript{h}akkaras}}, \textquotedblleft{}heap, pile, lump\textquotedblright{} — the flatbread named for the lump of dough it started as. That reconstruction is as far back as the sources trace it; its own root before Proto-Indo-Aryan is not established, so this lesson claims none rather than inventing one.

Unlike \textbf{\mr{पाणी}} and \textbf{\mr{दूध}}, \textbf{\mr{भाकरी}} is not a Sanskrit tatsama or tadbhava — it never touched classical Sanskrit. And unlike \textbf{\mr{चहा}}, it is no loanword either — a regional word, close kin to Gujarati \emph{bhākhrī} and Kannada \emph{bakkari}, and simply \textbf{not} Hindi's word for bread: Hindi asks for \textbf{\mr{रोटी}} \emph{roṭī}, an unrelated root. A Delhi shopkeeper would not recognize the request; a Pune one would hand you lunch.
\end{cousinweb}

\begin{culture}
Say the four requests once more and notice their genders: \textbf{\mr{पाणी}} neuter, \textbf{\mr{चहा}} masculine, \textbf{\mr{दूध}} neuter, \textbf{\mr{भाकरी}} feminine. Four everyday nouns, all three of Marathi's genders, in one short chapter — proof, alongside the verbs' present-tense endings, that the third gender Hindi lost is not a museum piece here.
\end{culture}

\subsection*{Your turn}
\begin{itemize}
\raggedright
  \item \emph{Say it:} all four requests, kṛpayā each time — \textbf{pāṇī, chahā, dūdh, bhākarī}
  \item \emph{Say it:} \textquotedblleft{}I eat bhakri\textquotedblright{} — \textbf{mī bhākarī khāto} or \textbf{khāte}, putting Chapter 7's \textbf{\mr{खाणे}} to work on a word it never had before
  \item \emph{Name:} which of the four is neither Sanskrit nor a loanword — \textbf{bhākarī}
\end{itemize}

\begin{tcolorbox}[breakable,colback=teal!4,colframe=teal!35!black,title={Before you move on}]
What did \textbf{\mr{भाकरी}}'s Proto-Indo-Aryan ancestor mean? (\textbf{A heap, a lump}.) What does Hindi ask for instead? (\textbf{\mr{रोटी}}, an unrelated root.) Give all four requests with their genders. (\textbf{\mr{पाणी}} n, \textbf{\mr{चहा}} m, \textbf{\mr{दूध}} n, \textbf{\mr{भाकरी}} f — all with \textbf{\mr{कृपया}}.) Sort the four sources: two from Sanskrit, one a Chinese loan, one purely regional. (\textbf{Pāṇī, dūdh} — Sanskrit; \textbf{chahā} — Chinese; \textbf{bhākarī} — regional.)

Sources: \href{https://en.wiktionary.org/wiki/\%E0\%A4\%AD\%E0\%A4\%BE\%E0\%A4\%95\%E0\%A4\%B0\%E0\%A5\%80}{Wiktionary: \mr{भाकरी}}.
\end{tcolorbox}
